\subsection{内乘：余代数情形}

设 \(E = \bigoplus_{p \geq 0} E_p\) 是一个余结合、余单位的分次余代数。那么我们知道（第 2 号，命题 1 和 3），分次对偶 \(E^{*gr}\)（按照本段开头关于分次的约定）具有 A 上的一个 N 型分次代数结构，该代数中两个元素 \(u, v\) 的乘积由 \(uv = m \circ (u \otimes v) \circ c\) 定义，其中 \(c: E \to E \otimes_A E\) 是余积，而 \(m: A \otimes_A A \to A\) 定义乘法。换言之，若对 \(x \in E\) 有 \(c(x) = \sum_i y_i \otimes z_i\) ，则我们可以写成（将 \(A \otimes_A E\) 与 \(E\) 典范地等同）

\[
\begin{array}{r l} \langle x, u v \rangle = (u v) (x) & = \sum_ {i} u (y _ {i}) v (z _ {i}) = v \left(\sum_ {i} u (y _ {i}) z _ {i}\right) \\ & = v (((u \otimes 1 _ {\mathbb {E}}) \circ c) (x)) = \langle ((u \otimes 1 _ {\mathbb {E}}) \circ c) (x), v \rangle . \end{array}
\]

这可以解释为：对 \(E^{*gr}\) 中所有次数为 p 的齐次元素 u，左乘 \(e(u): v \mapsto uv\) 在 \(E^{*gr}\) 中是 E 的次数为 -p 的分次自同态的分次转置

\[
i (u) = (u \otimes 1 _ {\mathrm{E}}) \circ c\tag{46}
\]

因此，在上述记号中，

\[
(i (u)) (x) = \sum_ {i} u (y _ {i}) z _ {i}.
\]

公式 (46) 同样定义了一个元素 \(i(u)\in\mathrm{Endgr}_{\mathbf{A}}(\mathbf{E})\)（对每个元素 \(u\in\mathbf{E}^{*\mathbf{g}\mathbf{r}}\) 而言）；对所有 \(x\in\mathbf{E}\) 及所有 \(u\in\mathbf{E}^{*\mathbf{g}\mathbf{r}}\) ，我们记

\[
x \llcorner u = (i (u)) (x)\tag{47}
\]

使得，对 E*gr 中的 u 和 v，

\[
\langle x, u v \rangle = \langle x \llcorner u, v \rangle .
\]

E 的元素 \(x \llcorner u\) 称为 x 被 u 的右内乘。

这里再次将“作用于”x的元素u放在L中水平线的自由端。

对任意两个元素 \(u, v\)（属于 \(\mathbf{E}^{*\mathrm{gr}}\)），

\[
x \llcorner (u v) = (x \llcorner u) \llcorner v,\tag{48}
\]

换句话说

\[
i (u v) = i (v) \circ i (u).
\]

如上所述，令 \(c(x) = \sum_{i} y_{i} \otimes z_{i}\) ，则有 \(x \llcorner (uv) = \sum_{i} (uv)(y_{i})z_{i}\) 。若

\[
c (y _ {i}) = \sum_ {j} y _ {i j} ^ {\prime} \otimes y _ {i j} ^ {\prime \prime},
\]

那么

\[
x \llcorner (u v) = \sum_ {i, j} u \left(y _ {i j} ^ {\prime}\right) v \left(y _ {i j} ^ {\prime \prime}\right) z _ {i}.\tag{49}
\]

另一方面，如果 \(c(z_{i}) = \sum_{k} z_{ik}' \otimes z_{ik}''\) ，则

\[
(x \llcorner u) \llcorner v = \sum_ {i, k} u (y _ {i}) v (z _ {i k} ^ {\prime}) z _ {i k} ^ {\prime \prime}.\tag{50}
\]

现在，E 的余结合性表明（第 2 号，命题 1）

\[
\sum_ {i, j} y _ {i j} ^ {\prime} \otimes y _ {i j} ^ {\prime \prime} \otimes z _ {i} = \sum_ {i, k} y _ {i} \otimes z _ {i k} ^ {\prime} \otimes z _ {i k} ^ {\prime \prime}\tag{51}
\]

而表达式 (49) 与 (50) 的相等性来自如下事实：它们分别是 (51) 式左、右两端在从 \(E \otimes E \otimes E\) 到 E 的线性映射 f 下的像，其中 \(f(x \otimes y \otimes z) = u(x)v(y)z\) 。

另一方面，我们回顾（第 2 号，命题 3），代数 \(\mathbf{E}^{*\mathrm{gr}}\) 的单位元是线性形式 \(e: x \mapsto \gamma(x).1\) ；因此

\[
x \llcorner e = \sum_ {i} \gamma (y _ {i}) z _ {i} = x
\]

这是余单位定义的直接推论。由于映射 \(u \mapsto i(u)\) 是线性的，可以看出外合成律 \((u, x) \mapsto x \llcorner u\) 在 E 上定义了一个右 E*gr-模结构。

类似地，我们对所有 \(u \in \mathrm{E}^{*\mathbf{g}\mathbf{r}}\) 定义 \(\mathbf{E}\) 的自同态

\[
i ^ {\prime} (u) = \left(1 _ {E} \otimes u\right) \circ c\tag{52}
\]

并且，对所有 \(x \in \mathbf{E}\) ，我们记

\[
\left(i ^ {\prime} (u)\right) (x) = u \lrcorner x\tag{53}
\]

并且 E 的这个元素称为 x 被 u 的左内乘。如前所述可以看出，外合成律 \((u, v) \mapsto u \lrcorner x\) 在 E 上定义了一个左 \(\mathrm{E}^{*gr}\)-模结构。此外，这两个结构是相容的，换言之，

\[
(u \lrcorner x) \llcorner v = u \lrcorner (x \llcorner v)\tag{54}
\]

这对 E*gr 中的 \(u, v\)（II，§ 1，第 14 号）成立。采用与上述相同的记号，(54) 的左边是 \(\sum_{i,j} u(z_i) v(y_{ij}') y_{ij}''\) ，右边是 \(\sum_{i,k} v(y_i) u(z_{ik}') z_{ik}''\) ；它们的相等性来自如下事实：它们分别是 (51) 式左、右两端在从 E ⊗ E ⊗ E 到 E 的线性映射 \(g\) 下的像，其中 \(g(x \otimes y \otimes z) = v(x) u(z)y\) 。

因此可以看出，E 上的两个外合成律在此集合上定义了一个 \((\mathrm{E}^{*}\mathrm{gr}, \mathrm{E}^{*}\mathrm{gr})\)-双模结构。

当余代数 E 余交换时，则 \(u \lrcorner x = x \llcorner u\) ，其中 \(x \in \mathbf{E}\) 和 \(u \in \mathrm{E}^{*\mathrm{gr}}\) 任取；当它余反交换（§ 4，第 9 号）且 \(u \in \mathrm{E}_p^*\) 、 \(x \in \mathrm{E}_q\) 时，我们可以写成 \(c(x) = \sum_{0 \leqslant j \leqslant q} \left( \sum_i y_{ij} \otimes z_{i,q-j} \right)\) ，其中 \(y_{ij}\) 和 \(z_{ij}\) 都属于 \(\mathrm{E}_j\) （对每个 \(j\) ），于是由假设

\[
\sum_ {i} z _ {i j} \otimes y _ {i, q - j} = (- 1) ^ {j (q - j)} \sum_ {i} y _ {i j} \otimes z _ {i, q - j}.
\]

由定义， \(x \llcorner u = \sum_{0 \leqslant j \leqslant q} \left( \sum_{i} u(y_{ij}) z_{i,q-j} \right)\) 和

\[
u \lrcorner x = \sum_ {0 \leqslant j \leqslant q} \left(\sum_ {i} u (z _ {i, q - j}) y _ {i j}\right).
\]

由于 \(u(y_{ij}) = 0\) （相应地 \(u(z_{i,q-j}) = 0\) ），除非 \(j = p\) （相应地 \(q - j = p\) ），由上述可知 \(u \lrcorner x = (-1)^{p(q-p)}x \llcorner u\) 。

最后，设 \(\phi: \mathrm{E} \to \mathrm{F}\) 是一个分次余代数态射；那么已经看到（第 2 号），分次转置 \(^t\phi: \mathrm{F}^{*\mathrm{gr}} \to \mathrm{E}^{*\mathrm{gr}}\) 是一个分次代数同态；因此，对 \(x \in \mathrm{E}, u, v\) （属于 \(\mathrm{F}^{*\mathrm{gr}}\) ），

\[
\begin{array}{r l} \langle \phi (x \llcorner {} ^ {t} \phi (u)), v \rangle & = \langle x \llcorner {} ^ {t} \phi (u), ^ {t} \phi (v) \rangle = \langle x, ^ {t} \phi (u) ^ {t} \phi (v) \rangle = \langle x, ^ {t} \phi (u v) \rangle \\ & = \langle \phi (x), u v \rangle = \langle \phi (x) \llcorner u, v \rangle \end{array}
\]

从而

\[
\phi (x) \llcorner u = \phi (x \llcorner {} ^ {t} \phi (u));\tag{55}
\]

并且类似地

\[
u \lrcorner \phi (x) = \phi (^ {t} \phi (u) \lrcorner x).\tag{56}
\]

换言之， \(\phi\) 是一个 \(\mathbf{F}^{*\mathrm{gr}}\) -同态，从 \((\mathrm{E}^{*\mathrm{gr}}, \mathrm{E}^{*\mathrm{gr}})\) -双模 E 到 \((\mathrm{F}^{*\mathrm{gr}}, \mathrm{F}^{*\mathrm{gr}})\) -双模 F，相对于环同态

\[
{ } ^ { t } \phi : \mathrm{F} ^ { * \mathrm{gr} } \rightarrow \mathrm{E} ^ { * \mathrm{gr} } .
\]

例子。特别地，当 E 是分次余代数 \(\mathsf{T}(\mathbf{M})\) , \(\mathsf{S}(\mathbf{M})\) 或 \(\bigwedge(\mathbf{M})\) 之一（M 为一个 A-模）（第 1 号，例 5、6 和 7）时，可以应用上述内容。为了显式求出如此得到的双模结构，我们再次将 \(\mathsf{T}(\mathbf{M})^{*\mathbf{g}\mathbf{r}}\) （相应地 \(\mathsf{S}(\mathbf{M})^{*\mathbf{g}\mathbf{r}}\) ，相应地 \(\bigwedge(\mathbf{M})^{*\mathbf{g}\mathbf{r}}\) ）中次数为 n 的齐次元素 f 等同于 \(M^{n}\) 上的一个 n-线性形式（相应地，对称 n-线性形式；相应地，交错 n-线性形式，也称为 n-形式）。只需就 M 的元素表出这些乘积

\[
\left(x _ {1} \otimes x _ {2} \otimes \dots \otimes x _ {p}\right) \llcorner f (\text { resp. } (x _ {1} x _ {2} \dots x _ {p}) \llcorner f,
\]

（相应地 \((x_{1} \wedge x_{2} \wedge \cdots \wedge x_{p}) \llcorner f)\) ），对 M 中元素的每一有限序列 \((x_{i})_{1 \leqslant i \leqslant p}\) ，以及左内乘的类似表达式。现在，定义 (46) 和 (52) 以及第 1 号的公式 (3)、(6) 和 (9) 分别给出：

\[
\left\{ \begin{array}{l} (x _ {1} \otimes x _ {2} \otimes \dots \otimes x _ {p}) \llcorner f = f \lrcorner (x _ {1} \otimes x _ {2} \otimes \dots \otimes x _ {p}) = 0 \\ (x _ {1} x _ {2} \ldots x _ {p}) \llcorner f = f \lrcorner (x _ {1} x _ {2} \ldots x _ {p}) = 0 \quad \text {for} \quad p <   n. \\ (x _ {1} \wedge x _ {2} \wedge \dots \wedge x _ {p}) \llcorner f = f \lrcorner (x _ {1} \wedge x _ {2} \wedge \dots \wedge x _ {p}) = 0 \end{array} \right.\tag{57}
\]

对 \(p \geqslant n\) ，我们分别有

\begin{enumerate}
\def\labelenumi{(\arabic{enumi})}
\setcounter{enumi}{57}
\tightlist
\item
\end{enumerate}

\[
\left(x _ {1} \otimes x _ {2} \otimes \dots \otimes x _ {p}\right) \llcorner f = f \left(x _ {1}, \dots , x _ {n}\right) x _ {n + 1} \otimes \dots \otimes x _ {p}\tag{59}
\]

\[
\left(x _ {1} x _ {2} \dots x _ {p}\right) \llcorner f = \sum_ {\sigma} f \left(x _ {\sigma (1)}, \dots , x _ {\sigma (n)}\right) x _ {\sigma (n + 1)} \dots x _ {\sigma (p)}\tag{60}
\]

\[
\left(x _ {1} \wedge x _ {2} \wedge \dots \wedge x _ {p}\right) \llcorner f = \sum_ {\sigma} \varepsilon_ {\sigma} f \left(x _ {\sigma (1)}, \dots , x _ {\sigma (n)}\right) x _ {\sigma (n + 1)} \wedge \dots \wedge x _ {\sigma (p)}
\]

（其中，在 (59) 和 (60) 中，求和是对排列 \(\sigma \in \mathfrak{S}_p\) 中在每个区间 \([1, n]\) 和 \([n + 1, p]\) （均为 \(\mathbf{N}\) 的区间）上递增的那些进行的）；并且类似地

\begin{enumerate}
\def\labelenumi{(\arabic{enumi})}
\setcounter{enumi}{60}
\tightlist
\item
\end{enumerate}

\[
f \lrcorner (x _ {1} \otimes x _ {2} \otimes \dots \otimes x _ {p}) = f (x _ {p - n + 1}, \dots , x _ {p}) x _ {1} \otimes x _ {2} \otimes \dots \otimes x _ {p - n}\tag{62}
\]

\[
f \lrcorner (x _ {1} x _ {2} \dots x _ {p}) = \sum_ {\sigma} f (x _ {\sigma (p - n + 1)}, \dots , x _ {\sigma (p)}) x _ {\sigma (1)} \dots x _ {\sigma (p - n)}\tag{63}
\]

\[
f \lrcorner (x _ {1} \wedge x _ {2} \wedge \dots \wedge x _ {p}) =
\]

\[
\sum_ {\sigma} \varepsilon_ {\sigma} f (x _ {\sigma (p - n + 1)}, \dots , x _ {\sigma (p)}) x _ {\sigma (1)} \wedge \dots \wedge x _ {\sigma (p - n)}
\]

（其中，在 (62) 和 (63) 中，求和是对排列 \(\sigma \in \mathfrak{S}_p\) 中在每个区间 \([1, p - n]\) 和 \([p - n + 1, p]\) （均为 \(\mathbf{N}\) 的区间）上递增的那些进行的）。

